\section{Introduction}
\label{sec:intro}

\IEEEPARstart{T}{he} Internet of Things (IoT) is increasingly used to monitor
people continuously: in the operating room, in the intensive care unit and
through wearable sensors. Much of the value of these deployments comes from
combining different kinds of data. In an operating room, monitors record
physiological waveforms while the hospital record holds the patient's
clinical history; in a wearable installation, inertial sensors record motion
while a camera records what the wearer sees. A model that sees both streams
can recognize states that neither reveals alone.

Such data cannot simply be collected in one place. Clinical records are
governed by consent and by data-protection law, and continuous multimodal
sensing produces volumes that are costly to move. Federated learning (FL)
\cite{ref1} addresses this by training a shared model where the data reside
and exchanging only model updates. Hierarchical FL \cite{ref2, ref3} adds an
edge tier between the clients and the cloud, so that most aggregation happens
close to the data and the wide-area link is used only once every several
rounds, which reduces both traffic and energy.

Keeping the data local, however, does not by itself keep it private. Model
updates are computed from the records, and an adversary can often infer from
a released model whether a particular record was used in training
\cite{ref14}. Differential privacy (DP) gives a formal guarantee against such
inference. Its standard training algorithm, differentially private stochastic
gradient descent (DP-SGD) \cite{ref6}, bounds (clips) the influence of each
training record and adds calibrated noise, so the guarantee protects whatever
unit is clipped. In a sensing deployment that unit is a short window of
signal, not a person, and every patient contributes many windows. A
record-level guarantee therefore answers whether a particular window was used,
whereas a clinical deployment must answer whether a particular
\emph{patient} was. Group privacy \cite{ref8} relates the two, but at the
group sizes found in our data it turns a nominal $\varepsilon = 1$ into
$\varepsilon \approx 8$ with $\delta \approx 0.09$, which supports no clinical
claim.

Existing work has addressed communication efficiency, privacy and multimodal
learning largely separately. The review in Section~\ref{sec:related}
identifies four gaps.
\begin{enumerate}
\item \emph{No patient-level guarantee across a hierarchy.} Hierarchical
designs \cite{ref3, ref4, ref5} reduce communication and energy but provide no
formal guarantee, and private designs \cite{ref9, ref10} provide one but do
not account for how an edge tier changes each client's privacy expenditure.
Recent systems join a hierarchy with privacy \cite{ref34, ref37} or privacy
with multimodal clinical data \cite{ref35, ref36}; to our knowledge, none
states its guarantee over an individual whose records are grouped before
clipping.
\item \emph{An unmeasured cost.} The utility cost of protecting individuals
rather than records has not been measured on multimodal data paired at the
source, nor related to the number of records each individual contributes.
\item \emph{An implicit condition.} Private training of a small head on frozen
pretrained features \cite{ref12, ref13} is valid end to end only if the
features were pretrained on a population disjoint from the protected one.
This condition is usually left implicit, and it is not known whether the
benefit comes from the pretraining or from privatizing fewer parameters.
\item \emph{Audits at the wrong unit.} Guarantees are usually reported from
mechanism parameters alone and, where audited, audited per record, which is
exactly what a record-level guarantee already protects.
\end{enumerate}

This paper proposes PP-HFM (Patient-Private Hierarchical Federated Multimodal
learning), a client--edge--cloud framework in which the privacy unit is the
patient by construction. Each client sums the gradients of a patient's records
before clipping, so that the patient as a whole has bounded influence, and
calibrates its noise to its own patient population. The main contributions
are as follows.
\begin{itemize}
\item \emph{Patient-level privacy with complete accounting (gap 1).} A
patient-level private training procedure for hierarchical multimodal FL, with
per-client noise calibration, a proof of its guarantee
(Theorem~\ref{thm:guarantee}), and a ledger of every release, edge tier
included, that states which configurations are covered end to end.
\item \emph{A law for the cost of the unit (gap 2).} A scaling hypothesis,
motivated by excess-risk bounds (Proposition~\ref{prop:scaling}), under which
the utility gap between record- and patient-level training grows as
$A(1 - 1/k)$ with the number of records per individual $k$ and vanishes at
$k = 1$ (Corollary~\ref{cor:k1}), tested out of sample and on an external
corpus.
\item \emph{What makes a strict budget affordable (gap 3).} A three-arm design
showing that pretrained initialization, not the smaller privatized subset,
accounts for robustness at small $\varepsilon$.
\item \emph{System validation and an individual-level audit (gaps 1--4).} An
evaluation on three multimodal corpora paired at the source, with up to $187$
hospital clients; uplink traffic measured from the deployed encoding, delivery
and round latency simulated at packet level in ns-3, energy modeled in
CloudSim~Plus and inference latency timed on a workstation CPU; a patient-level membership-inference audit with a positive control;
and a comparison with re-implemented HED-FL and FL2DP under identical data,
encoders and seeds.
\end{itemize}
PP-HFM builds on user-level differential privacy \cite{ref11} and
parameter-efficient private fine-tuning \cite{ref12, ref13}; the
contributions are the system that composes them across all tiers, its proof,
and its evaluation on three multimodal corpora.

The rest of the paper is organized as follows. Section~\ref{sec:related}
reviews related work. Section~\ref{sec:model} states the problem, the system
and threat models and the research objectives. Section~\ref{sec:method}
presents PP-HFM and its privacy analysis. Section~\ref{sec:setup} describes the
experimental setup, Section~\ref{sec:results} reports the results,
Section~\ref{sec:discussion} discusses implications and limitations, and
Section~\ref{sec:conclusion} concludes with directions for future work.
